\chapter{Dry Socket}
\index{Dry Socket}

\section*{Definition and Overview}
Dry socket, known medically as alveolar osteitis, is a painful condition that can develop after a tooth is removed, most commonly after the extraction of wisdom teeth. Normally, a blood clot forms in the empty socket to protect the underlying bone and nerve endings while healing takes place. If the clot is lost, dissolves, or fails to form, the bone and nerve endings are exposed, causing severe, throbbing pain. Dry socket is not an infection, although it can become one, and it is a common complication that is easily treated by a dentist.

\section*{Epidemiology}
Dry socket is one of the most common complications of tooth extraction, occurring in a small but significant proportion of extractions. It is most frequent after the removal of impacted lower wisdom teeth and is more common in women, in smokers, and in people with poor oral hygiene. The condition usually develops within the first few days after the extraction, typically between the second and fifth day, when the blood clot is lost.

\section*{Aetiology and Pathophysiology}
After a tooth is removed, the socket fills with blood, which clots to protect the site. Dry socket occurs when this clot is lost or fails to form, exposing the bone and the alveolar nerve endings. The exposed bone becomes inflamed, and pain follows.
\begin{itemize}
    \item \textbf{Loss of the blood clot:} vigorous rinsing, spitting, drinking through a straw, or smoking in the first days after extraction can dislodge the clot
    \item \textbf{Smoking:} tobacco smoke irritates the socket and reduces blood supply, and the sucking action of smoking can dislodge the clot
    \item \textbf{Poor oral hygiene:} bacteria in the mouth can break down the clot
    \item \textbf{Difficult extractions:} traumatic or prolonged extractions increase the risk
    \item \textbf{Oral contraceptives:} higher oestrogen levels may affect clot formation
    \item \textbf{Past history:} people who have had a dry socket before are more likely to develop one again
\end{itemize}

\section*{Clinical Features}
Symptoms typically begin two to five days after the extraction.
\begin{itemize}
    \item Severe, throbbing pain in the socket that radiates to the ear, eye, or neck on the same side
    \item A partially or completely empty socket, visible where the blood clot should be
    \item Exposed bone visible in the socket
    \item A bad taste in the mouth or foul odour
    \item Slight swelling and redness of the gums around the socket
    \item Pain that is not relieved by over-the-counter pain relievers
\end{itemize}

\section*{Differential Diagnosis}
Other causes of pain after tooth extraction should be considered.
\begin{itemize}
    \item Normal post-extraction pain, which usually peaks within the first 24 hours and then settles
    \item Infection of the socket, which may cause swelling, pus, and fever
    \item A retained tooth fragment or foreign body in the socket
    \item Referred pain from other teeth, the jaw joint, or ear conditions
    \item Bruising or haematoma formation
\end{itemize}

\section*{When to Seek Medical Care}
Anyone who develops severe pain in the days after a tooth extraction should return to their dentist promptly. Dry socket is easy to diagnose and treat, and delaying treatment prolongs the pain. Call the dentist or the after-hours dental service if pain is severe or worsening.

\textbf{Contact your local emergency services for:}
\begin{itemize}
    \item Swelling spreading to the face, neck, or eye, which may indicate a serious infection
    \item Difficulty breathing or swallowing
    \item Fever with severe swelling
\end{itemize}

\section*{Diagnosis and Investigations}
Diagnosis of dry socket is clinical.
\begin{itemize}
    \item \textbf{History:} recent tooth extraction and the pattern of pain
    \item \textbf{Examination:} inspection of the socket to confirm that the blood clot is missing and bone is exposed
    \item \textbf{Gentle probing:} may confirm exposed bone and tenderness
    \item \textbf{Imaging:} dental X-rays are sometimes used to exclude retained fragments or other problems
\end{itemize}

\section*{Management}
Dry socket is treated in the dental surgery, and the pain usually resolves quickly with treatment.
\begin{itemize}
    \item \textbf{Socket irrigation:} the dentist gently cleans the socket with saline to remove debris
    \item \textbf{Medicated dressing:} a dressing containing a pain-relieving and antiseptic agent is placed in the socket, often providing rapid relief
    \item \textbf{Dressing changes:} the dressing may need to be replaced every few days until healing progresses
    \item \textbf{Pain relief:} over-the-counter analgesics such as paracetamol and ibuprofen, and sometimes stronger medication
    \item \textbf{Oral hygiene:} gentle rinsing with warm salt water once the socket begins to heal
    \item \textbf{Antibiotics:} only if there is evidence of infection
\end{itemize}

\section*{Complications}
\begin{itemize}
    \item Prolonged severe pain, which can interfere with eating, sleeping, and daily activities
    \item Infection of the socket, which may spread to the jawbone in rare cases
    \item Delayed healing of the extraction site
    \item Temporary numbness or altered sensation in the lip, chin, or tongue in rare cases
\end{itemize}

\section*{Prognosis and Outcomes}
Dry socket is a temporary condition that resolves completely once the socket begins to heal. With treatment, most people experience significant pain relief within a day or two, and the socket heals over the following weeks. The condition does not cause long-term damage to the jaw, and future extractions can be managed carefully to reduce the risk of recurrence.

\section*{Prevention and Self-Care}
\begin{itemize}
    \item Follow the dentist's aftercare instructions carefully after an extraction
    \item Avoid vigorous rinsing, spitting, and drinking through a straw for the first few days
    \item Do not smoke for at least 48 to 72 hours after the extraction, and longer if possible
    \item Maintain gentle oral hygiene, keeping the socket clean without disturbing it
    \item Avoid hot or crunchy foods that could dislodge the clot
    \item Rest and avoid strenuous activity for the first day
\end{itemize}

\section*{Sources and Further Reading}
The following reputable sources provide further information for healthcare students and patients seeking reliable, current advice about dry socket.
\begin{itemize}
    \item NHS. \textit{Dry socket: symptoms and treatment.} https://www.nhs.uk/conditions/dry-socket/
    \item Mayo Clinic. \textit{Dry socket: symptoms, causes, and treatment.} https://www.mayoclinic.org/diseases-conditions/dry-socket/
    \item MedlinePlus. \textit{Dry socket after tooth extraction.} https://medlineplus.gov/ency/article/001044.htm
    \item Australian Dental Association. \textit{After an extraction: patient advice.} https://www.ada.org.au/
    \item Cleveland Clinic. \textit{Dry socket: causes and prevention.} https://my.clevelandclinic.org/health/diseases/dry-socket
    \item MSD Manual. \textit{Alveolar osteitis: professional overview.} https://www.msdmanuals.com/
\end{itemize}
